% Generated from tex/chapters/ch04/07_実務上の考慮.tex for the SWEBOK structure tree
\subsubsection{構築テスト}

構築では、主に単体テストと統合テストが行われる。これらは、多くの場合、コードを書いたソフトウェア技術者自身が実行する。構築テストの目的は、欠陥が混入してから発見されるまでの時間差を短くし、修正費用を下げることである。

単体テストは、関数、クラス、モジュールなどを単独で確認する。統合テストは、それらを組み合わせたときに正しく相互作用するかを確認する。構築テストは、通常、システムテスト、アルファテスト、ベータテスト、ストレステスト、使用性テストなどの広い範囲のテストまでは含まない。

テストケースは、コードを書いた後に作られることもあるが、コードを書く前に作られることもある。テストを先に書く方法は、テストファーストプログラミングまたはテスト駆動開発として扱われる。
